\chapter{Código Fuente y Datos Analizados}

Este capítulo reúne la trazabilidad de los artefactos técnicos del trabajo: \emph{cómo} se accede o reproduce el código y \emph{qué} conjuntos de datos soportan el análisis, sin sustituir el detalle de la metodología, ya recogida en el Capítulo 3.

\section{Estructura del repositorio y módulos principales}
Describa la organización del repositorio (carpetas, módulos que implementan VAE, GAN y la optimización de cartera) y, si aplica, el flujo de ejecución (scripts de entrenamiento, evaluación, generación de figuras o tablas). Puede acompañarse de un esquema textual o de referencia a nombres de archivos; evite embeber listados extensos de código en el PDF si su normativa UNIR aconseja anexar el código o depositarlo en repositorio público.

\section{Fuentes de datos y criterios de muestreo}
Indique el origen de las series o paneles (mercado, frecuencia, ventana temporal), criterios de limpieza y alineación, y el tratamiento de ausencias o cotizaciones. Justifique la elección de activos, benchmark o universo, en coherencia con el diseño de experimentos.

\section{Disponibilidad y condiciones de uso}
Consigne dónde reside el material (repositorio, almacenamiento, anexos) y las condiciones o licencias que afecten a datos o terceros. Si el código no es público, indíquese la forma aceptada de entregarlo a evaluación según le indique su tutoría.

\section{Reproducibilidad}
Especifíquense variables de entorno, \textit{seeds} para resultados estocásticos, o ficheros de configuración que permitan reproducir los resultados clave, en la medida en que la confidencialidad o las restricciones de los datos lo permitan.
